% resumo em inglês
\setlength{\absparsep}{18pt} % ajusta o espaçamento dos parágrafos do resumo
\begin{resumo}[Abstract]

\begin{otherlanguage*}{english}
   
Discussions on entrepreneurship in the Amazon region highlight the importance of local brands in sustainable actions, but little attention has been paid to brand management strategies, especially those of small companies. Decision-making activities that use new methods in environmental areas often involve extensive research in abundant literature. Such activities are laborious, time-consuming, and difficult to be interpreted by entrepreneurs. From this perspective, it seems advantageous to develop safe, automated, organized and controllable mechanisms for collecting evidence on management models. This mechanism should make it easier for users to optimize the time and work required to maintain the reliability of the results. In this sense, this research aims to formulate a theoretical-methodological model for validating brands based on elements related to environmental sustainability in the Brazilian Amazon. The research consists of understanding two fronts of information, which involves the perception of the brand equity of 8 (eight) entrepreneurs and their corresponding brands, as well as the perception of potential consumers of these brands guided by the concept of sustainability based on six formative constructs: Willingness to buy WTB; Brand strategic positioning — SBP; perception of product quality — BQP and Innovation in retail — IR Brand Authenticity — BA. The results are presented in five blocks: first, a Cochrane scope review was carried out, seeking to explore studies on branding and sustainable brands, as well as detailing the latest advances in environmental entrepreneurship for the agrarian environment with extensive activity, by reviewing the scope of the literature international based on the understanding of the main existing dimensions and scales, a scale was developed and its psychometric properties of accuracy and validity validated, for assessing brand equity, based on environmental sustainability. The scale was validated by adopting the standards described in the Norm of standards for educational and psychological testing. Next, an identification was made of which Brand Equity constructs based on brands related to the environmental sustainability of the Amazon influence consumer loyalty to the brand and willingness to purchase its products. In the fourth moment, semi-structured interviews were carried out using the Grounded Theory, aiming to understand the perception of entrepreneurs from their business model related to the connection of the Amazon brand. The fifth moment of the research proposes to evaluate the effects of three dimensions that make up Brand Equity in the authenticity of brands related to the environmental sustainability of the Amazon. From all the findings, the last moment aimed to prototype a theoretical-methodological model that uses the environmental sustainable Brand Equity of the Amazon of Brazil for the development of new sustainable brands.

\textbf{Keywords}: Authenticity of Agroforestry Brands; Original Forest Communities; Commercialization of Agroforestry Products; Amazon of Brazil.
 \end{otherlanguage*}
\end{resumo}